% Appendix 6: Equations of Planes in Three Dimensions
% VIC Specialist Mathematics — Units 3–4; QLD Specialist Mathematics — Unit 3
\begingroup
\small

\noindent\fbox{%
\begin{minipage}{0.97\textwidth}
\textit{\textbf{Beyond the HSC syllabus.}
The NESA Mathematics Extension~2 syllabus requires understanding of
coordinate planes ($z = 0$, etc.)\ and coplanarity through vector
relationships---see \S\ref{sec:planes-hsc-note} in the Vectors
Primer.
This appendix covers the \textbf{full algebra of general plane
equations} as taught in
\textbf{VIC Specialist Mathematics (Units~3--4)} and
\textbf{QLD Specialist Mathematics (Unit~3)}.
It is intended as enrichment for university preparation,
\textbf{not} as additional HSC examination content.
Students should master the coordinate-plane reasoning in the main
booklet before exploring this material.}
\end{minipage}}

\vspace{0.8em}

% ---- Key concepts ----
\noindent\textbf{Key concepts.}

In three-dimensional space, a line can be determined by a point and a
direction vector:
\[
  \mathbf{r} = \mathbf{p} + t\,\mathbf{d}.
\]
A \textbf{plane} can be determined by a point and a non-zero
\textbf{normal vector}.

A normal vector is perpendicular to every direction lying in the
plane.

\vspace{0.4em}
\noindent\textbf{Parametric form of a plane.}
A plane can also be described using a point and \textbf{two
non-parallel direction vectors} that lie in the plane:
\[
  \boxed{
  \mathbf{r} = \mathbf{p} + s\,\mathbf{u} + t\,\mathbf{v}
  },
  \qquad s, t \in \mathbb{R}.
\]
This is analogous to the parametric form of a line
$\mathbf{r} = \mathbf{a} + \lambda\,\mathbf{b}$,
but with \emph{two} independent parameters instead of one.

\vspace{0.4em}
\noindent\textbf{From parametric form to Cartesian form.}
Given two direction vectors $\mathbf{u}$ and $\mathbf{v}$ in the
plane, a normal vector can be found using the cross product
(Appendix~\ref{app:cross-product}):
\[
  \mathbf{n} = \mathbf{u} \times \mathbf{v}.
\]
Then the point-normal equation is
\[
  \boxed{\mathbf{n} \cdot (\mathbf{r} - \mathbf{p}) = 0}.
\]
If $\mathbf{n} = (a, b, c)$ and
$\mathbf{p} = (x_0, y_0, z_0)$, this becomes
\[
  a(x - x_0) + b(y - y_0) + c(z - z_0) = 0.
\]
Equivalently, the Cartesian equation of a plane is
\[
  \boxed{ax + by + cz = d},
\]
where $a$, $b$, $c$ are not all zero.

% ---- Worked example ----
\vspace{0.8em}
\noindent\textbf{Worked example --- Finding a plane from three
points.}

\noindent Find the Cartesian equation of the plane passing through
\[
  A(1, 0, 0), \qquad B(2, 0, 1), \qquad C(1, 2, 0).
\]

\vspace{0.3em}
\noindent\textit{Step~1: Find two vectors in the plane.}
\[
  \overrightarrow{AB} = (1, 0, 1),
  \qquad
  \overrightarrow{AC} = (0, 2, 0).
\]

\noindent\textit{Step~2: Find a normal vector using the cross
product.}

From Appendix~\ref{app:cross-product},
\[
  \overrightarrow{AB} \times \overrightarrow{AC}
  = \bigl(0\cdot 0 - 1\cdot 2,\;
         1\cdot 0 - 1\cdot 0,\;
         1\cdot 2 - 0\cdot 0\bigr)
  = (-2,\, 0,\, 2).
\]
A simpler normal vector is therefore
\[
  \mathbf{n} = (1, 0, -1).
\]

\noindent\textit{Step~3: Form the plane equation.}

Using point $A(1, 0, 0)$,
\[
  (1, 0, -1) \cdot (x - 1,\; y,\; z) = 0.
\]
Expanding,
\[
  x - 1 - z = 0.
\]
Hence,
\[
  \boxed{x - z = 1}.
\]

\noindent\textit{Verification.}
All three given points satisfy the equation:
$A$: $1 - 0 = 1$~$\checkmark$;\quad
$B$: $2 - 1 = 1$~$\checkmark$;\quad
$C$: $1 - 0 = 1$~$\checkmark$.

% ---- Key takeaways ----
\vspace{0.8em}
\noindent\textbf{Key takeaways.}
\begin{itemize}[leftmargin=*,itemsep=2pt]
  \item A line is described by one direction vector; a plane can be
        described by two non-parallel direction vectors.
  \item The cross product provides a quick method of finding a
        plane's normal vector.
  \item Coordinate planes such as $z = 0$ are special cases of
        general plane equations.
\end{itemize}

\vspace{0.4em}
\noindent\textit{Beyond school --- University connection.}
Plane equations appear in linear algebra, computer-aided design,
3D graphics and optimisation.
They also form the foundation for studying intersections of surfaces
in multivariable calculus.

\endgroup
